\begin{figure}[htbp]
\centering
\begin{tikzpicture}[
 node distance=9mm and 13mm,
 box/.style={draw=black!65,rounded corners=2pt,align=center,
             text width=39mm,minimum height=10mm,font=\small,inner sep=4pt},
 arrow/.style={-{Stealth[length=2mm]},semithick}]
 \node[box] (trans) {Transverse power saving};
 \node[box,below=of trans] (sparse) {Sparse packet estimate};
 \node[box,right=of trans] (frame) {Uniform bilinear frame estimate};
 \node[box,below=of frame] (flat) {Flat-frequency estimates};
 \node[box,right=of flat] (geom) {Rich-tube geometry};
 \node[box,below=14mm of flat] (dist) {Uniform dyadic\\distribution bound};
 \node[box,below=of dist] (sum) {Endpoint square summation};
 \node[box,below=of sum] (limit) {Gaussian limit at all times};
 \draw[arrow] (trans)--(sparse);
 \draw[arrow] (frame)--(flat);
 \draw[arrow] (frame.east) -- ++(6mm,0) |- ([yshift=2.5mm]dist.east);
 \draw[arrow] (sparse.south)|-(dist.west);
 \draw[arrow] (flat)--(dist);
 \draw[arrow] (geom.south)|-(dist.east);
 \draw[arrow] (dist)--(sum);
 \draw[arrow] (sum)--(limit);
\end{tikzpicture}
\caption{Dependencies of the proof. Small scale losses are allowed in the
transverse and sparse arguments; the frame and dyadic distribution estimates
must be uniform in the packet scale.}
\label{fig:dependencies}
\end{figure}
